\section{Анализ существующих работ}

\subsection{Внутренняя калибровка и оценка позы}

Метод Zhang~\cite{zhang2000} использует несколько изображений плоского шаблона для определения внутренних параметров камеры и последующего нелинейного уточнения. Для рассматриваемой задачи такая калибровка является предварительным этапом: параметры проекции должны быть известны либо оценены отдельно. Калибровка отдельной камеры сама по себе не задаёт взаимное положение всех камер установки.

При известных внутренних параметрах положение объекта относительно камеры можно оценивать по соответствиям между его трёхмерными точками и их проекциями. В работе Lepetit и соавторов~\cite{lepetit2009} предложен метод EPnP, вычислительная сложность которого линейна по числу соответствий. Это подходящая основа для получения начальной оценки позы, однако качество результата зависит от распределения наблюдаемых точек и точности их локализации. В прикладной реализации дополнительно необходимы отбраковка ошибочных соответствий и проверка геометрической допустимости решения.

\subsection{Маркеры и многокамерная калибровка}

Система AprilTag~\cite{olson2011} позволяет обнаруживать и идентифицировать искусственные визуальные метки. Идентификаторы помогают сопоставлять точки изображения с заранее заданной геометрией объекта. Однако наличие меток не определяет автоматически координаты их углов в системе всей мишени: расположение меток на гранях должно быть задано и проверено отдельно. Выбор семейства маркеров в настоящей работе будет зависеть от фактической конструкции калибровочного объекта.

Детектирование шаблона в ToF-изображениях требует отдельного внимания. Hansard и соавторы~\cite{hansard2014} рассматривают обнаружение шахматной сетки с использованием её глобальной геометрии. Эта работа показывает значимость специализированной обработки таких изображений; она не даёт оснований заранее считать произвольный детектор маркеров работоспособным на данных ещё не изученной камеры.

Hansard и соавторы~\cite{hansard2015} рассматривают взаимную калибровку нескольких систем ToF и цветных камер с использованием калибровочного шаблона. Эта работа непосредственно относится к объединению разных наблюдений в общую геометрическую модель. Для нашей постановки существенны доступность изображения мишени, известность её геометрии и согласованность систем координат каналов. При этом применение конкретной схемы измерений должно определяться доступными данными установки.

Совместное уточнение параметров камер и структуры наблюдений относится к классу задач bundle adjustment, рассмотренному Triggs и соавторами~\cite{triggs2000}. Здесь структура мишени предполагается известной, а оптимизируемыми величинами служат внешние параметры камер и позы объекта в отдельных наблюдениях. Такая постановка позволяет использовать все наблюдения в одной функции ошибки и явно зафиксировать систему отсчёта.

\subsection{Особенности глубины ToF и регистрация}

Обзор Foix и соавторов~\cite{foix2011} систематизирует принципы работы и калибровки lock-in ToF-камер, а также источники ошибок измерения. Из этого следует необходимость учитывать качество глубины при её включении в геометрическую оптимизацию. В настоящей постановке измерения около границ объекта и недостоверные отсчёты предлагается исключать, а оставшиеся остатки нормировать с учётом принятой модели погрешности.

Другой подход к совмещению измерений --- регистрация трёхмерных данных. Алгоритм ICP Besl и McKay~\cite{besl1992} итеративно уточняет преобразование на основе соответствий между геометрическими элементами. Применительно к нескольким камерам такой подход требует достаточной общей геометрии и подходящего начального приближения. Если камеры видят разные грани мишени, прямое сопоставление облаков может оказаться недостаточно определённым. Поэтому в качестве базового метода здесь выбрана привязка к известной геометрии объекта, а регистрация облаков рассматривается как возможный дополнительный способ проверки или уточнения при наличии перекрытия.

Таким образом, рассматриваемая схема опирается на известные методы, а исследовательский вопрос состоит в их применимости к доступной ToF-установке, устойчивости к неполным наблюдениям и чувствительности к погрешностям мишени. Преимущество использования глубины по сравнению с калибровкой только по изображениям должно быть проверено экспериментально.
